\chapter{Rapid Tests}
\section*{What Is This Test?} Rapid tests provide preliminary results near the patient, often within minutes, using a specimen such as blood, urine, saliva, or a swab.\section*{Alternative Names} Point-of-care testing; bedside testing; near-patient testing.\section*{Principle} Immunoassay, molecular, chemical, or electrode methods detect an analyte or pathogen quickly.\section*{Why This Test Is Ordered} Rapid tests support immediate triage, treatment, infection control, glucose monitoring, pregnancy testing, or screening.\section*{Primary vs Secondary Test} A rapid result may be primary or preliminary; confirmatory laboratory testing depends on the test and clinical context.\section*{How This Test Is Performed} Trained staff or patients collect the specified specimen and follow the device instructions and quality controls.\section*{What Preparation Is Needed} Test-specific instructions apply; timing, collection, storage, and operator technique affect accuracy.\section*{Understanding Results} Positive and negative predictive value varies with prevalence; a negative rapid test may require confirmation when suspicion is high.\section*{Follow-up Tests} Laboratory confirmation, culture, PCR, repeat testing, or clinical evaluation may follow.\section*{References} \begin{enumerate}\item MedlinePlus. Rapid Tests.\item U.S. Food and Drug Administration. Home and point-of-care testing information.\end{enumerate}\clearpage\section*{Understanding Results and Follow-up}
The laboratory report should identify the specimen, method, units, reference interval or decision limit, and whether the result is positive, negative, abnormal, or inconclusive. Reference intervals vary with age, sex, pregnancy, specimen type, assay, and laboratory; a result outside a range is not automatically a diagnosis, and a result within a range does not always exclude disease. Discuss unexpected results with the ordering clinician, who can decide whether repeat or confirmatory testing, treatment, or referral is appropriate. Do not change medicines or delay urgent care based on an isolated result.
\section*{Regulatory Status and Authorization Date}This chapter describes a test category, not a specific commercial product. FDA, CDSCO, EU IVDR, and MHRA approval or registration dates therefore cannot be assigned without the assay manufacturer, product name, instrument, and intended use. For laboratory-developed tests, an individual agency approval date may not exist. See the Preface for the interpretation of regulatory status.
\clearpage
